\documentclass[journal,12pt, onecolumn]{IEEEtran}

% Packages

\usepackage{cite}
\usepackage{graphicx}
\usepackage{mathtools}
\usepackage[warnings-off={mathtools-colon,mathtools-overbracket}]{unicode-math}
\usepackage{braket}
\usepackage{hyperref}
\usepackage{fontspec}

\usepackage{readable-source}

\setmainfont{Cambria}
\setmathfont{Cambria Math}

% Initializations
\hbadness=99999
\DeclareGraphicsExtensions{.pdf,.jpeg,.png}



\title{Review of HSP for Deutsch's algorithm}
\author{
Aarohi Bhavsar \\
Aarohi.Bhavsar@rockets.utoledo.edu
Date: November 12, 2026 \\
Supervisor: Hossam Shalaby
The University of Toledo
}

\begin{document}
    \maketitle
    \begin{abstract}
        \par{
            Deutsch's algorithm is the one of the earliest instances of solving a case of the Hidden Subgroup. This algorithm efficiently solves for a bit of partial information, $c$, about the definition of a function
        }
        
    \end{abstract}
    \eq{f:G\rightarrow{}X,}
        
        given
        \begin{equation}
             G=\B{0,1},\quad X=\B{0,1},\quad (\text{and its quantum gate formulation}, U_f)
        \end{equation}
                   
    
        
    \begin{figure}[h]
        \centering
        \includegraphics[width=1\textwidth]{Deutsch's Circuit.png}
        \caption{Deutsch's Circuit}
    \end{figure} 
    \section{State Analysis}

    \begin{equation}
        \begin{aligned}
            |𝜓_0 ⟩ &= |0⟩_{𝑞_0} |1⟩_{𝑞_1}
            \\
            |𝜓_1 ⟩& = (𝐻⊗𝐻)|𝜓_0 ⟩
            \\
            &= (𝐻⊗𝐻) |0⟩_{𝑞_0} |1⟩_{𝑞_1}
            \\
            &= |+⟩_{𝑞_0} |−⟩_{𝑞_1}
            &\qquad &= ½(|0⟩_{𝑞_0} |0⟩_{𝑞_1}−|0⟩_{𝑞_0} |1⟩_{𝑞_1}+|1⟩_{𝑞_0} |0⟩_{𝑞_1}−|1⟩_{𝑞_0} |1⟩_{𝑞_1})
            \\
            |𝜓_2 ⟩ &= 𝑈_𝑓 |𝜓_1 ⟩
            \\
            &= 𝑈_𝑓 |+⟩_{𝑞_0} |−⟩_{𝑞_1}
            &\qquad &= ½(𝑈_𝑓 |0⟩_{𝑞_0} |0⟩_{𝑞_1}−𝑈_𝑓 |0⟩_{𝑞_0} |1⟩_{𝑞_1}+𝑈_𝑓 |1⟩_{𝑞_0} |0⟩_{𝑞_1}−𝑈_𝑓 |1⟩_{𝑞_0} |1⟩_{𝑞_1})
        \end{aligned}
    \end{equation}

    
    % \eqaligned{
    %     {\forall}C, {C\prime}:
    %     \p{{\exists}g,g\prime{}\in{}G, C} &= g{\circ}H
    %     {\wedge}
    %         {C\prime} = {g\prime}{\circ}H\\&
    %         f\p{C} = \B{f\p{c}:c{\in}C}
    %     {\wedge}
    %         f\p{C} = f\p{C\prime}{\leftrightarrow}{C\prime}= C{\prime}
    %     {\wedge}
    %         {\forall}c,{c\prime}{\in}C,f\p{c}=f\p{c{\prime}}
    % }
        
    
\end{document}


